\section{Reachable sections at arbitrary width}\label{sec:lifts51}
In this section $D\ge2$, $X=\{\pm e_1,\ldots,\pm e_D\}$, all
coordinate queries are present, $I\in\calA\subset O(D)$, and
$0<\rho<1$. A command cut is indexed by $t\in\{0,\ldots,N\}$;
$T_{t+1,a}$ sends cut $t$ to cut $t+1$. All statements concern the
clocked, horizon-specific model of Section~\ref{sec:machine}.
Write $\Delta_{k-1}=\{p\in\R^k:p\ge0,\ p\mathbf1=1\}$.
Distributions and observable means are row vectors in matrix identities;
physical vectors are their transposes.

The correct observable object above minimum rank is a section followed
by a projection. This is the reachable-space construction underlying
positive realization, not a license to identify every hidden basis
state with a physical state. We give its controlled form, including the
number of generators needed after projection.

\begin{theorem}[Reachable-section geometry]\label{thm:section51}
Let an exact realization have profile $(K_t)$. Let $E_t$ have as rows
all actual state distributions after prefixes $(x,u)$ of length $t$,
and let $Z_t\in[-1,1]^{K_t\times D}$ have as columns the conditional
means of the coordinate queries after the identity suffix. Set
\[
 R_t=\operatorname{rowspan}E_t,\quad r_t=\dim R_t,\quad
 P_t=R_t\cap\Delta_{K_t-1},\quad S_t=\{(pZ_t)^T:p\in P_t\}.
\]
Then $r_t\ge D+1$, $\dim P_t=r_t-1$, and $S_t$ is a
full-dimensional polytope satisfying
\begin{equation}\label{eq:section51}
 C_\rho\subset S_t\subset Q_t\subset[-1,1]^D,\qquad
 U_aS_t\subset S_{t+1},\qquad
 f_0(S_t)\le {K_t\choose r_t-1}.
\end{equation}
Here $C_\rho=\rho\conv X$ and
$Q_t=\{z:|e_j^TU_vz|\le1\text{ for }|v|=N-t,\ 1\le j\le D\}$.
Moreover $P_tT_{t+1,a}\subset P_{t+1}$, and for $s<t$,
$U_wS_s\subset S_t$ for every word of length $t-s$, without any
restriction on intervening widths. The affine fibers of the map
$p\mapsto pZ_t$ on $R_t\cap\{p\mathbf1=1\}$ have dimension
$r_t-D-1$.
\end{theorem}
\begin{proof}
For a suffix $v$, denote by $B_t(v)\in[-1,1]^{K_t\times D}$ its
actual conditional mean matrix. Exactness, first on actual prefix rows
and then on their linear span, gives
\begin{equation}\label{eq:span51}
 yB_t(v)=yZ_tU_v^T\quad(y\in R_t).
\end{equation}
In particular this identity is not asserted for $y$ outside $R_t$.
If $p\in P_t$, the left side is a convex combination of legal mean
rows. Hence $(pZ_t)^T\in Q_t$. Identity prefixes give points
$\pm\rho e_i$ in $S_t$, and identity suffixes give its cube bound.
The map $y\mapsto(y\mathbf1,yZ_t)$ has rank $D+1$ on $R_t$:
its values on the identity-prefix seed rows are $(1,\pm\rho e_i^T)$,
which span $\R^{D+1}$. Thus $r_t\ge D+1$ and the stated fiber
dimension follows from rank--nullity.

The affine hull of the reachable rows is precisely
$R_t\cap\{p\mathbf1=1\}$, since each of those rows has mass one.
Their convex hull lies in $P_t$ and has affine dimension $r_t-1$;
therefore $P_t$ has this dimension as well. On this affine space the
only inequalities defining $P_t$ are its $K_t$ nonnegativity
inequalities. At each vertex some $r_t-1$ active coordinate
inequalities have linearly independent restrictions to the translation
space. Such a choice uniquely determines the vertex. Assigning one
choice to each vertex injects the vertex set into the subsets of size
$r_t-1$. Projection cannot increase the number of vertices, proving
the binomial bound. Redundant inequalities, including permanently zero
coordinates, cause no difficulty in the independence argument.

For each actual prefix row $p$, $pT_{t+1,a}$ is a row of $E_{t+1}$.
Linearity gives $R_tT_{t+1,a}\subset R_{t+1}$, and stochasticity gives
the inclusion of sections. Exactness with identity padding gives
\[
 yT_{t+1,a}Z_{t+1}=yZ_tU_a^T\quad(y\in R_t).
\]
This proves the physical containment; composition proves its
separated-cut form.
\end{proof}

\begin{corollary}[One surplus label]\label{cor:surplus51}
At a cut of width $D+2$ there are exactly two possible reachability
ranks. If $r_t=D+2$, then $P_t=\Delta_{D+1}$ and $S_t$ has at most
$D+2$ vertices, with a one-dimensional affine observation kernel.
If $r_t=D+1$, then $P_t$ is a $D$-dimensional section of
$\Delta_{D+1}$ and its observation map is an affine isomorphism.
Consequently $S_t$ has at most $D+2$ facets and at most
$\binom{D+2}{2}$ vertices. At width $D+1$ the section is the whole
simplex and the observation map is an affine isomorphism.
\end{corollary}
\begin{proof}
Theorem~\ref{thm:section51} leaves only the two indicated values of
$r_t$. Its dimension and kernel statements give each assertion. An
affine isomorphism preserves facets; the section has at most $D+2$
coordinate facets. At width $D+1$ the kernel has dimension zero.
\end{proof}

This distinguishes two forms of surplus: unused linear directions in
the ambient register, of codimension $K_t-r_t$, and unobservable
directions within the reachable span, of dimension $r_t-D-1$.
The image $S_t$ alone does not retain the stochastic extension
requirement. In particular, counting its vertices does not identify
its static extension complexity with the original state width.

\subsection{A local polynomial characterization}
The complete suffix-response arrays are useful for interpretation but
are not required as input to an exact feasibility formula.

\begin{theorem}[Bounded local realization criterion]\label{thm:local51}
For prescribed positive integers $(K_t)$, an exact realization exists
if and only if the following finite system has a real solution.
For each cut choose a symmetric orthogonal projection
$\Pi_t\in\R^{K_t\times K_t}$ and a matrix
$Z_t\in[-1,1]^{K_t\times D}$; for each letter choose a row-stochastic
$T_{t+1,a}\in\R^{K_t\times K_{t+1}}$; and choose
$p_x\in\Delta_{K_0-1}$ for every $x\in X$. Require
\begin{align}
 \Pi_t^T&=\Pi_t,\qquad \Pi_t^2=\Pi_t,\label{eq:projector51}\\
 p_x\Pi_0&=p_x,\qquad p_xZ_0=\rho x^T,\label{eq:seed51}\\
 \Pi_tT_{t+1,a}(I-\Pi_{t+1})&=0,\label{eq:invariant51}\\
 \Pi_t(T_{t+1,a}Z_{t+1}-Z_tU_a^T)&=0.\label{eq:intertwine51}
\end{align}
The variables may all be restricted to $[-1,1]$. The system has degree
at most three and at most
\begin{equation}\label{eq:variables51}
 |X|K_0+\sum_{t=0}^N(K_t^2+DK_t)
       +|\calA|\sum_{t=0}^{N-1}K_tK_{t+1}
\end{equation}
scalar variables. Its description is polynomial in the explicitly
expanded clock length, the profile, the alphabet size and $D$, and
contains no enumeration of command words. For algebraic input it is
a finite decision and synthesis problem over the real algebraic
numbers.
\end{theorem}
\begin{proof}
For necessity use the actual $R_t$ and $Z_t$ of
Theorem~\ref{thm:section51} and the orthogonal projector onto $R_t$.
The span identities just proved give
\eqref{eq:invariant51}--\eqref{eq:intertwine51}; initialization gives
\eqref{eq:seed51}. An orthogonal projector has operator norm at most
one, and hence entries of absolute value at most one. All other
variables have the asserted bounds by stochasticity and legality.

Conversely let $p$ be a running distribution with $p\Pi_t=p$.
Equation~\eqref{eq:invariant51} gives
$(pT_{t+1,a})\Pi_{t+1}=pT_{t+1,a}$, while
\eqref{eq:intertwine51} gives
$pT_{t+1,a}Z_{t+1}=pZ_tU_a^T$. Starting with $p_x$, induction
therefore yields the exact terminal mean $\rho(U_wx)^T$. Use the
columns of $Z_N$ as decoder means. Each is bounded by one, so it
defines a legal binary decoder on every available state, including
states never reached. This constructs the required common-row machine.

The equations have degree at most three, and counting their displayed
matrices gives \eqref{eq:variables51}. All equalities and weak
inequalities are closed in a compact box. For algebraic input,
real-algebraic elimination and sampling decide nonemptiness and produce
an algebraic point when nonempty \cite{BPR}. The criterion is valid for
arbitrary real data as a mathematical equivalence; effective input
encoding in the last assertion is a separate assumption.
\end{proof}

The bound concerns the number of variables and equations, not the
running time of quantifier elimination. In particular $N$ is the
number of listed epochs, not its binary encoding length. No
polynomial-time higher-width optimizer is inferred.

\begin{proposition}[Positivity-compatible local alternative]\label{prop:dual51}
Fix the projections and mean matrices in
Theorem~\ref{thm:local51}. For one transition, stack its row-sum and
intertwining equations as $B\operatorname{vec}(T)=b$. A stochastic
extension exists if and only if
\[
 B^Ty\ge0\quad\Longrightarrow\quad b^Ty\ge0
 \qquad\hbox{for every real }y.
\]
An infeasible transition has a finite witness $B^Ty\ge0$, $b^Ty<0$.
For rational supplied data this is an exact rational linear-programming
alternative.
\end{proposition}
\begin{proof}
The remaining inequalities are exactly
$\operatorname{vec}(T)\ge0$. Apply the equality form of Farkas'
lemma. The row-sum equations are included in $B$, so nonnegativity
also bounds every entry by one. Rational polyhedral feasibility and
separation give rational witnesses when the data are rational.
\end{proof}

The local alternative is classical linear separation, applied to the
full positivity constraint. It is not a global convex dual over
unknown sections. Its necessity can be seen without any dimension
asymptotics. On
\[
 Q=\{p\in\Delta_3:p_1+p_2=p_3+p_4=1/2\},
 \qquad F(p)=(2p_1,2p_2),
\]
$F$ is affine and maps $Q$ into $\Delta_1$, but its best ambient
stochastic extension has uniform total-variation error $1/4$. Indeed,
write $u_i=T_{i1}\in[0,1]$. At
$(e_1+e_j)/2,(e_2+e_j)/2$, $j=3,4$, the two desired probabilities
differ by one whereas the realized probabilities differ by at most
$1/2$. One error is at least $1/4$. The choice
$u=(1,0,1/2,1/2)$ attains this bound at all four vertices, and
convexity bounds the error on their hull. Thus affine positivity on a
section cannot replace \eqref{eq:invariant51}--\eqref{eq:intertwine51}
with an ambient nonnegative transition.
